\chapter{Các công trình liên quan}
\label{Chapter2}

\textit{Chương này điểm lại các công trình liên quan làm nền tảng cho khóa luận: các kiến trúc học sâu cho phân đoạn ảnh y sinh (U-Net và các biến thể), những hướng tiếp cận của học máy lượng tử cùng các mô hình lai lượng tử -- cổ điển, từ đó chỉ ra những khoảng trống nghiên cứu mà khóa luận hướng tới giải quyết.}

\section{Học sâu cho phân đoạn ảnh y sinh}
Trong những năm gần đây, học sâu (Deep Learning) \cite{lecun2015deeplearning} đã tạo ra một cuộc cách mạng trong lĩnh vực phân tích ảnh y sinh, đặc biệt là bài toán phân đoạn ảnh. Kể từ thành công của các mạng nơ-ron tích chập sâu trong phân loại ảnh \cite{krizhevsky2012alexnet, simonyan2015vgg, he2016resnet}, hướng tiếp cận học sâu cho phân đoạn ảnh cũng nhanh chóng phát triển: mạng tích chập toàn phần (Fully Convolutional Network) \cite{long2015fcn} lần đầu cho phép dự đoán nhãn ở mức điểm ảnh đầu-cuối, tiếp nối bởi các kiến trúc mã hóa -- giải mã như SegNet \cite{badrinarayanan2017segnet} và DeepLab \cite{chen2018deeplabv3plus}. Trong bối cảnh đó, kiến trúc U-Net do Ronneberger và cộng sự đề xuất vào năm 2015 \cite{ronneberger2015unet} được coi là cột mốc quan trọng nhất cho phân đoạn ảnh y sinh. Với cấu trúc đối xứng dạng chữ U bao gồm một nhánh mã hóa (encoder) để trích xuất đặc trưng không gian và một nhánh giải mã (decoder) để phục hồi độ phân giải, kết hợp cùng các kết nối tắt (skip connections) giúp bảo tồn chi tiết vị trí, U-Net đã giải quyết xuất sắc bài toán phân đoạn với số lượng dữ liệu huấn luyện hạn chế.

Sự thành công của U-Net đã thúc đẩy sự ra đời của hàng loạt các biến thể cải tiến. UNet++ \cite{zhou2018unetplusplus} đưa ra kiến trúc các mạng con kết nối lồng nhau (nested skip pathways) nhằm thu hẹp khoảng cách ngữ nghĩa (semantic gap) giữa encoder và decoder, còn UNet 3+ \cite{huang2020unet3plus} mở rộng thành kết nối tắt toàn tỷ lệ (full-scale). Attention U-Net \cite{oktay2018attention} tích hợp cơ chế tập trung (attention mechanism) vào các kết nối tắt, giúp mô hình tự động chọn lọc vùng cần phân tích và ức chế các đặc trưng nhiễu từ nền ảnh; R2U-Net \cite{alom2018r2unet} bổ sung các khối tích chập hồi quy -- phần dư. Đối với dữ liệu khối ba chiều, 3D U-Net \cite{cicek20163dunet} và V-Net \cite{milletari2016vnet} tổng quát hóa kiến trúc lên không gian thể tích, trong khi nnU-Net \cite{isensee2021nnunet} chứng minh rằng một quy trình tự cấu hình dựa trên U-Net vẫn là đường cơ sở (baseline) rất mạnh trên nhiều thử thách y sinh. Gần đây hơn, việc du nhập cơ chế self-attention của Transformer \cite{vaswani2017attention, dosovitskiy2021vit} đã sinh ra các kiến trúc lai như TransUNet \cite{chen2021transunet}, Swin-Unet \cite{cao2022swinunet} và Medical Transformer \cite{valanarasu2021medt}, tận dụng khả năng mô hình hóa phụ thuộc tầm xa (long-range dependencies).

Một hướng cải tiến khác, không dựa trên attention mà vẫn mở rộng được trường nhìn (receptive field) với chi phí thấp, là \emph{tích chập giãn nở} (dilated/atrous convolution) \cite{yu2016dilated}: bằng cách chèn các ``lỗ hổng'' vào bộ lọc, một nhân tích chập $3\times 3$ có thể bao phủ vùng ngữ cảnh $5\times 5$ hay $7\times 7$ mà không tăng số tham số. Ý tưởng tổng hợp ngữ cảnh đa tỷ lệ này được hệ thống hóa trong khối Atrous Spatial Pyramid Pooling của DeepLab \cite{chen2018deeplabv3plus}, và kết hợp với tích chập tách kênh theo chiều sâu (depthwise separable convolution) \cite{chollet2017xception} để giảm mạnh chi phí tính toán. Đây chính là nền tảng cho khối tích chập đa tỷ lệ (RDMS) mà khóa luận sử dụng để củng cố bộ khung cổ điển của mô hình lai, sẽ được trình bày ở Chương~\ref{Chapter4}. Tuy nhiên, dù có bao nhiêu biến thể được phát triển, kiến trúc cơ sở cốt lõi của U-Net vẫn thường được xem là nền tảng chuẩn mực và vững chắc cho các tác vụ phân đoạn y tế. Hạn chế chung của tất cả các phương pháp này là chúng hoàn toàn dựa trên điện toán cổ điển, vốn gặp khó khăn trong việc biểu diễn các hàm phi tuyến độ phức tạp cực cao khi kiến trúc bị giới hạn về độ sâu và số lượng tham số.

\section{Học máy lượng tử trong xử lý ảnh y tế}
Sự ra đời của Học máy Lượng tử (Quantum Machine Learning - QML) đã mở ra một hướng đi mới để vượt qua các giới hạn của điện toán cổ điển \cite{biamonte2017qml, schuld2015introqml}. Thông qua các Mạch lượng tử tham số hóa (PQC), QML hứa hẹn cung cấp không gian trạng thái Hilbert có số chiều tăng trưởng theo hàm mũ, mang lại khả năng biểu diễn (expressivity) vô cùng phong phú \cite{du2025qmltutorial, schuld2019featurehilbert}. Nhiều mô hình lượng tử tham số hóa đã được đề xuất, từ quantum circuit learning \cite{mitarai2018qcl} và mạng nơ-ron lượng tử cho thiết bị gần hạn \cite{farhi2018qnn}, tới việc xem PQC như những mô hình học máy tổng quát \cite{benedetti2019pqc} và mạng tích chập lượng tử QCNN \cite{cong2019qcnn}. Havl{\'i}{\v{c}}ek và cộng sự đã chứng minh thực nghiệm lợi thế của không gian đặc trưng lượng tử trên phần cứng siêu dẫn \cite{havlicek2019supervised}, trong khi các khảo sát về thuật toán biến phân \cite{cerezo2021vqa} và về năng lực của mạng nơ-ron lượng tử \cite{abbas2021power} đặt nền tảng lý thuyết cho hướng đi này. Việc huấn luyện các mô hình lượng tử được hỗ trợ bởi các khung vi phân tự động như PennyLane \cite{bergholm2018pennylane}, song vẫn phải đối mặt với thách thức ``cao nguyên cằn cỗi'' khi mạch quá sâu \cite{mcclean2018barren}, nhất là trong điều kiện phần cứng NISQ hiện tại \cite{preskill2018nisq}.

Trong lĩnh vực ảnh y khoa, các công trình kết hợp mạng nơ-ron học sâu với điện toán lượng tử (Hybrid Quantum-Classical Neural Networks - HQCNN) đang thu hút sự chú ý mạnh mẽ. Tuy nhiên, một thực tế dễ nhận thấy là phần lớn các nghiên cứu hiện tại chỉ tập trung giải quyết bài toán phân loại (classification). Ví dụ, Shahjalal và cộng sự (2025) \cite{shahjalal2025hqcnn} đã đề xuất mô hình HQCNN cho tác vụ phân loại ảnh y khoa và đạt được những kết quả khả quan về độ chính xác và khả năng hội tụ nhanh trên các tập dữ liệu nhỏ.

Đáng chú ý, các ứng dụng QML dành riêng cho bài toán phân đoạn ảnh (segmentation) vẫn còn cực kỳ thưa thớt do tính phức tạp của việc khôi phục cấu trúc không gian (spatial reconstruction) ở đầu ra thay vì chỉ dự đoán một nhãn duy nhất. Gần đây nhất, kiến trúc QC-Net (2025) \cite{long2025qcq} và khối trích xuất đặc trưng lượng tử QuFeX của Jain và Kalev (2026) \cite{qufex2025} đã bước đầu khám phá khả năng thay thế hoặc kết hợp các lớp tích chập cổ điển bằng các bộ trích xuất lượng tử. Đây chính là công trình mà khóa luận chọn làm điểm tham chiếu chính để phát triển và so sánh (mô hình đối chứng "QuNet" ở Chương~\ref{Chapter5}), nên đáng nói chi tiết hơn: QuFeX gom nhiều bản đồ đặc trưng đầu vào thành từng nhóm nhỏ để xử lý song song qua \emph{cùng một} mạch lượng tử dùng chung tham số, nhờ đó chỉ cần $O(1)$ lượt gọi mạch bất kể số bản đồ đặc trưng -- giảm mạnh chi phí so với xử lý riêng lẻ từng bản đồ như các hướng tiếp cận QCNN trước đó. Khối này được nhúng vào đúng vùng nút thắt cổ chai của U-Net (kiến trúc lai gọi là Qu-Net), bao quanh bởi một kết nối phần dư giữ nguyên ánh xạ đồng nhất. Trên chính tập PH2 mà khóa luận cũng sử dụng, nhóm tác giả báo cáo Qu-Net cải thiện IoU khoảng $10$ điểm phần trăm so với U-Net cổ điển, đồng thời phương sai giữa các lần chạy thấp hơn hẳn -- dù cần lưu ý đây là số liệu do chính nhóm tác giả công bố, theo quy trình đánh giá riêng của họ, không nhất thiết cùng thang đo hay điều kiện huấn luyện với thực nghiệm của khóa luận. Khóa luận kế thừa ý tưởng cốt lõi của công trình này -- đặt một khối tính toán lượng tử tại vùng nút thắt cổ chai của U-Net, bao bọc bởi kết nối phần dư -- làm điểm khởi đầu tham chiếu, nhưng thay khối lượng tử \emph{đơn} bằng một cơ chế \emph{hỗn hợp nhiều nhánh xử lý lượng tử đa tỷ lệ}. Dù vậy, những công trình trên vẫn chưa đưa ra được một đánh giá có hệ thống về việc kết hợp cơ chế nhiều nhánh xử lý lượng tử với các khối mô hình hóa chuỗi hiệu quả tại vùng nút thắt cổ chai của U-Net. Đây chính là khoảng trống nghiên cứu (research gap) mà khóa luận này hướng đến. Để làm rõ giá trị của cơ chế đề xuất, mô hình lai được so sánh trong điều kiện kiểm soát nghiêm ngặt với một phổ rộng các mô hình cổ điển tiêu biểu -- U-Net \cite{ronneberger2015unet}, UNet++ \cite{zhou2018unetplusplus}, Attention U-Net \cite{oktay2018attention}, V-Net \cite{milletari2016vnet} và R2U-Net \cite{alom2018r2unet} -- vốn có số tham số lớn hơn nhiều lần.

\section{Mixture of Experts trong Học sâu}
Ý tưởng Mixture of Experts (MoE) có nguồn gốc từ công trình kinh điển của Jacobs và cộng sự (1991) về hỗn hợp các chuyên gia cục bộ thích nghi \cite{jacobs1991moe}, sau đó được đưa vào các mạng học sâu nhiều tầng \cite{eigen2013moe}. Khái niệm này thực sự tạo tiếng vang lớn trong cộng đồng học sâu hiện đại thông qua lớp MoE định tuyến thưa của Shazeer và cộng sự (2017) \cite{shazeer2017moe}, và tiếp tục được mở rộng ở quy mô rất lớn với Switch Transformer \cite{fedus2022switch} cùng các mô hình thị giác như V-MoE \cite{riquelme2021vmoe}. Về bản chất, MoE là một khung kiến trúc (framework) bao gồm nhiều mạng con song song, thường được gọi là các ``chuyên gia'' (experts) trong tài liệu gốc -- khóa luận gọi các mạng con này là \emph{nhánh xử lý} để tránh gây hiểu nhầm rằng chúng có vai trò chuyên môn hóa như con người -- đi kèm với một bộ định tuyến (gating network / router). Với cơ chế định tuyến linh hoạt (ví dụ: softmax gating), mỗi đầu vào hoặc mỗi vùng không gian của ảnh sẽ được chuyển hướng tới một vài nhánh xử lý phù hợp nhất, cho phép tăng đáng kể số lượng tham số và năng lực của mô hình mà không làm tăng tuyến tính chi phí tính toán thực tế.

Mặc dù MoE đã mang lại những đột phá lớn cho các mô hình ngôn ngữ lớn (LLMs) và mô hình thị giác máy tính cổ điển (Vision Transformers), việc đưa cơ chế này vào không gian tính toán lượng tử chỉ vừa mới bắt đầu được khám phá trong thời gian rất gần đây. Ở cấp độ khung tổng quát, Nguyen và cộng sự (2025) đề xuất QMoE \cite{nguyen2025qmoe} -- một khung Quantum Mixture of Experts trong đó mỗi nhánh xử lý là một mạch lượng tử tham số hóa cùng một bộ định tuyến lượng tử học được -- nhưng được khảo sát chủ yếu cho tác vụ \emph{phân loại}. Gần với bài toán của khóa luận hơn, HQF-Net (2026) \cite{hossain2026hqfnet} lần đầu đưa một khối ``Quantum bottleneck with Mixture-of-Experts (QMoE)'' vào bài toán \emph{phân đoạn ảnh}, kết hợp các mạch lượng tử cục bộ, toàn cục và có hướng thông qua một cơ chế định tuyến thích nghi -- tuy nhiên trên dữ liệu \emph{viễn thám} (remote sensing).

Như vậy, ứng dụng Quantum MoE cho \emph{ảnh y sinh} vẫn còn là một khoảng trống chưa được khai phá, đặt trong bối cảnh dữ liệu y tế có những đặc thù rất khác (kích thước nhỏ, ranh giới mờ, mất cân bằng lớp). Quan trọng hơn, điểm mới về mặt \emph{cơ chế} của khóa luận không nằm ở việc ``dùng MoE lượng tử'', mà ở một thiết kế cụ thể chưa từng xuất hiện: các nhánh xử lý lượng tử được phân hóa theo \emph{kích thước mảnh ảnh} (patch size) -- tạo ra tính đa tỷ lệ ngay trong không gian lượng tử -- và được \emph{phối hợp với một khối mô hình hóa không gian trạng thái lấy cảm hứng từ Mamba} (LSS) tại cùng vùng nút thắt cổ chai để bổ khuyết ngữ cảnh toàn cục. Sự kết hợp đồng thời ``Quantum MoE đa tỷ lệ theo patch $+$ khối chuỗi lấy cảm hứng từ Mamba (LSS) $+$ khối tích chập đa tỷ lệ RDMS'' cho phân đoạn ảnh y sinh, theo hiểu biết của nhóm, chưa từng được đề xuất trong bất kỳ công trình nào trước đó. Thông qua cơ chế này, mô hình được kỳ vọng sẽ tự động điều phối (route) các vùng ảnh có cấu trúc từ đơn giản đến phức tạp vào mạch lượng tử phù hợp nhất, giải quyết bài toán trích xuất đặc trưng không gian ở vùng nút thắt cổ chai (bottleneck).

\section{Mô hình không gian trạng thái và Mamba}
Một hướng nghiên cứu song song và bùng nổ gần đây là mô hình không gian trạng thái (State Space Model - SSM) cho mô hình hóa chuỗi. Khởi đầu từ kiến trúc S4 \cite{gu2022s4}, các SSM có cấu trúc cho phép nắm bắt phụ thuộc tầm xa với độ phức tạp tuyến tính theo chiều dài chuỗi, khắc phục chi phí bậc hai của cơ chế self-attention. Kiến trúc Mamba \cite{gu2023mamba} tiếp tục cải tiến bằng cơ chế quét chọn lọc (selective scan), cho phép các tham số của mô hình phụ thuộc vào đầu vào, qua đó chọn lọc thông tin cần ghi nhớ một cách linh hoạt tương tự attention nhưng vẫn giữ chi phí tuyến tính; mối liên hệ hình thức giữa Mamba và cơ chế attention sau đó được làm rõ trong Mamba-2 \cite{dao2024mamba2}. Nhờ ưu thế này, Mamba nhanh chóng được ứng dụng sang lĩnh vực thị giác thông qua các kiến trúc như Vision Mamba \cite{zhu2024visionmamba} và VMamba \cite{liu2024vmamba}, và đặc biệt vào phân đoạn ảnh y sinh với U-Mamba \cite{ma2024umamba} cùng VM-UNet \cite{ruan2024vmunet}, như một giải pháp thay thế nhẹ và hiệu quả cho Transformer. Khóa luận kế thừa ý tưởng này ở dạng một khối chuỗi rút gọn -- khối Lightweight Selective Scan (LSS) -- đặt sau khối QuantumMoE, nhằm bổ khuyết khả năng tổng hợp thông tin toàn cục vốn còn hạn chế khi các nhánh xử lý lượng tử chỉ xử lý theo từng mảnh ảnh cục bộ.

\section{Khoảng trống nghiên cứu và hướng tiếp cận của khóa luận}
Bốn hướng nghiên cứu được điểm qua ở trên -- kiến trúc U-Net và biến thể cho phân đoạn y sinh, học máy lượng tử lai, cơ chế Mixture of Experts, và mô hình không gian trạng thái Mamba -- cho đến nay phần lớn vẫn phát triển tương đối độc lập. Các mô hình lai lượng tử phần lớn tập trung vào bài toán phân loại, và các nỗ lực đưa Quantum MoE vào phân đoạn mới chỉ xuất hiện rất gần đây trên dữ liệu ngoài y sinh \cite{hossain2026hqfnet}; cơ chế Mixture of Experts được khai thác chủ yếu trong không gian cổ điển; còn Mamba mới chỉ được kết hợp với các thành phần thuần cổ điển. Khóa luận này tiếp cận ở giao điểm của cả bốn hướng: sử dụng U-Net làm bộ khung, thay thế vùng nút thắt cổ chai bằng một cơ chế hỗn hợp nhiều nhánh xử lý \emph{lượng tử} đa tỷ lệ, bổ sung một khối \emph{Mamba} (LSS) để tổng hợp phụ thuộc không gian tầm xa, và củng cố bộ khung cổ điển bằng khối tích chập đa tỷ lệ \emph{RDMS}. Việc kết hợp đồng thời các thành phần này trong một kiến trúc phân đoạn ảnh y sinh -- được đặt tên là \textbf{HQMoSS-Net} -- theo hiểu biết của nhóm, chưa được khảo sát trong các công trình trước đó, và đó chính là đóng góp trọng tâm của khóa luận sẽ được trình bày chi tiết ở Chương~\ref{Chapter4}.